% ------------------------------------------------------------------------
% ------------------------------------------------------------------------
% ------------------------------------------------------------------------
%                                Resumen
% ------------------------------------------------------------------------
% ------------------------------------------------------------------------
% ------------------------------------------------------------------------
\chapter*{RESUMEN}

\footnotesize{
\begin{description}
  \item[TÍTULO:] \MakeUppercase{\titulo}
  \astfootnote{Trabajo de grado}
  \item[AUTOR:]{\autora}, {\autorb}
  \asttfootnote{\facultad. \escuela. Director: \director. Codirectores: \codirectora}
  \item[PALABRAS CLAVE:] RETINOPATÍA DIABÉTICA, REDES NEURONALES CONVOLUCIONALES, APRENDIZAJE PROFUNDO, IMÁGENES DEL FONDO DEL OJO, PROCESAMIENTO DIGITAL DE IMÁGENES, APRENDIZAJE POR TRANSFERENCIA.
  \item[DESCRIPCIÓN:] \hfill \\La retinopatía diabética (RD) es una complicación microvascular de la diabetes que daña los vasos sanguíneos de la retina, afectando a 145 millones de personas a nivel mundial y constituyendo una de las principales causas de ceguera prevenible. En Colombia, con una proporción de 1.56 oftalmólogos por cada 100,000 habitantes y aproximadamente 1.2 millones de personas con retinopatía diabética, la detección automatizada representa una solución crítica para ampliar la cobertura del tamizaje. El diagnóstico convencional mediante exámenes de fondo de ojo resulta costoso y limitado por la disponibilidad de especialistas, especialmente en regiones con recursos médicos escasos. Este trabajo desarrolla modelos basados en redes neuronales convolucionales (CNN) para detectar la enfermedad y clasificar sus niveles de severidad mediante clasificadores binarios (presencia/ausencia) y multiclase (cinco niveles). Se utilizó aprendizaje por transferencia con las arquitecturas MobileNetV2, ResNet50V2, InceptionV3 y VGG19, entrenadas con 35,121 imágenes del dataset público de Kaggle. El preprocesamiento incluyó normalización de color, enmascaramiento periférico, balanceo de clases mediante submuestreo y aumento de datos, incorporando bloques Multi-Head Attention, capas GaussianNoise y técnicas de regularización (Dropout, L2, early stopping). El modelo fue evaluado mediante sensibilidad, precisión, F1-score y AUC-ROC, siendo MobileNetV2 con preprocesamiento mediante filtrado y umbral optimizado (0.4826) el modelo de mejor desempeño, alcanzando una sensibilidad de 96.5\% en clasificación binaria con intervalos de confianza bootstrap (IC95\%: AUC-ROC 0.9636–0.9936). Para clasificación multiclase, se alcanzó una exactitud de 55\% y sensibilidad de 65\%, identificando el canal verde como el más efectivo debido a su mejor contraste para lesiones hemorrágicas. La validación externa en datasets MESSIDOR-2, IDRiD y FOSCAL confirmó la capacidad de generalización del modelo, demostrando su viabilidad como herramienta de apoyo al tamizaje clínico de retinopatía diabética.
\end{description}}\normalsize
% ------------------------------------------------------------------------
